\documentclass[../../lecture-notes-148x210.tex]{subfiles}

\begin{document}

This section introduces a new class of zero-knowledge protocols which are based
on the Sum-Check protocol. If previously we mostly worked over univariate
polynomials, Sum-Check derivatives typically operate over multilinear
polynomials in many variables. Turns out that such representation can give a
competitive performance to protocols relying on high-degree univariate
polynomials. 

Besides Sum-Check-based approaches, we also introduce the notion of
\textit{lookup checks}, which is a quite recent development, allowing to
significantly reduce the arithmetical circuit size for non-native problems for
the a zero-knowledge protocol (such as hash function computation or
zero-knowledge Machine Learning). While this is a standalone research direction,
we still provide it after discussing Sum-Check protocol since historically one
of the first lookup protocols (namely, logup) is in fact Sum-Check-based.

Specifically, we list sections of this chapter in \Cref{tab:essentialls-4}. 

\begin{table}[H]
    \centering
    \begin{tabularx}{\textwidth}{ccX}
        \Xhline{3\arrayrulewidth}
        \rowcolor{yellow!30}\textbf{Section} & \textbf{Topic} & \textbf{Key Concepts} \\
        \Xhline{3\arrayrulewidth}
        \rowcolor{yellow!10}\ref{section:sumcheck} & Sum-Check Protocol & Construction of 
        Sum-Check protocol, its performance and soundness  \\
        \hdashline
        \rowcolor{yellow!20}\ref{section:gkr} & GKR Protocol &
        Sum-Check-based SNARK over arithmetical circuits \\
        \hdashline
        \rowcolor{yellow!10}\ref{section:lookups} & Lookup Checks & Motivation for lookup checks; plookup and logup constructions \\
        \hdashline
        \rowcolor{yellow!20}\ref{section:ultragroth} & UltraGroth & Implementing lookup checks in Groth16 \\
        \Xhline{3\arrayrulewidth}
    \end{tabularx}
    \caption{Topics covered in Part IV}
    \label{tab:essentialls-4}
\end{table}

\end{document}